% Generated draft: requires mathematical and source review.
\documentclass[11pt]{article}
\usepackage[margin=25mm]{geometry}
\usepackage{fontspec}
\setmainfont{Latin Modern Roman}
\usepackage{amsmath,amssymb,mathtools,unicode-math}
\setmathfont{Latin Modern Math}
\usepackage{enumitem,needspace,xurl,hyperref}
\setlength{\parindent}{0pt}
\setlength{\parskip}{5pt}
\setlength{\emergencystretch}{3em}
\newcommand{\N}{\mathbb{N}}
\newcommand{\Z}{\mathbb{Z}}
\newcommand{\Q}{\mathbb{Q}}
\newcommand{\R}{\mathbb{R}}
\newcommand{\C}{\mathbb{C}}
\newcommand{\F}{\mathbb{F}}
\newcommand{\A}{\mathbb{A}}
\newcommand{\PP}{\mathbb{P}}
\newcommand{\E}{\mathbb{E}}
\newcommand{\eps}{\varepsilon}
\newcommand{\dd}{\,\mathrm{d}}
\begin{document}
\section*{UnsolvedMath ID: 20002297}
\textbf{Source problem code:} catalogue entry
\subsection{Definitions and mathematical statement}
Vidaux's Question 13 asks for applications of exceptional sets in higher-power Büchi problems in characteristic \(p\). It is a one-line research direction: no ground field, power exponent, sequence length, definition of exceptional set, or desired application is specified. The aid reports that Pasten--Wang classify appropriate positive-characteristic higher-power exceptions as pseudo-powers and apply them to existential definability and undecidability, while expressly not asserting an unrestricted solution in every setting. Their precise hypotheses and definition of pseudo-power are not reproduced.

The aid's explicit finite-subfield obstruction has the following hypotheses. Let \(p\) be prime, \(n\ge2\) with \(\gcd(p,n)=1\), \(d\ge1\), and let \(K\) be a characteristic-\(p\) field containing a specified constant subfield \(\F_{p^d}\). Assume \(K\) has an element \(f\notin K^p=\{u^p:u\in K\}\). Choose a positive integer \(E\) with
\[
 d\mid E,\qquad p^E\equiv1\pmod n,
\]
and set
\[
 F(X)=(X-f)^{n-1}(X-f^{p^E})\in K[X].
\]
The reported conclusions are that \(F\) is monic of degree \(n\), is a Frobenius pseudo-power in the aid's terminology, is not an actual \(n\)-th power in \(K[X]\), yet
\[
 F(b)\in K^n=\{u^n:u\in K\}
 \quad\text{for every }b\in\F_{p^d}.
\]
The value identity making the reported claim explicit is
\[
 F(b)=(b-f)^{n-1+p^E}
 =\left((b-f)^{(n-1+p^E)/n}\right)^n.
\]
The coprimality permits the choice of \(E\), and \(f\notin K^p\) excludes collision of the two roots. Their multiplicities \(n-1\) and one distinguish the polynomial from an actual \(n\)-th power. The restriction \(n\ge2\) is an explicit higher-power qualification; at \(n=1\) a non-first-power conclusion would be impossible.

Thus for every prescribed finite constant subfield there is a member of this exceptional family passing every power-value test on that subfield. The member and \(E\) may depend on \(d\); no single polynomial is asserted to pass simultaneously over all finite constant extensions. The statement concerns values that are powers in \(K\), not necessarily powers inside the finite test subfield. It obstructs eliminating this exception by any one prescribed finite-subfield test, without deciding the full unspecified Büchi or definability program.
\subsection{Short English statement}
The source asks broadly how higher-power Büchi exceptions in positive characteristic can be used. The aid reports literature applications and an explicit Frobenius family showing that every prescribed finite constant subfield can be saturated by power values without the polynomial itself being a power. The field, exponent, and order of quantifiers are made explicit.
\subsection{Sources}
\begin{itemize}
\item Local UnsolvedMath registry, permanent ID 20002297, exact source code catalogue entry; canonical \texttt{aim-logic-notes.json}, \texttt{notes[72]}, tag \texttt{question}; AIM Workshop Problem Lists, Logic, \emph{Definability and Decidability Problems in Number Theory}, Question 13 (Vidaux), 2013 according to supplied verification. Created and updated 14 August 2026.
\item Source URL \url{https://aimath.org/pastworkshops/definabilityinntproblems.pdf}. The supplied record is the complete one-line question. The PDF was not accessed.
\item Pasten--Wang is named for classification of higher-power positive-characteristic pseudo-power exceptions and applications to existential definability and undecidability. No full title, publication date, theorem number, separate URL, or complete hypothesis list is supplied.
\item Attempt 1: \texttt{partial\_result}, problem status at run \texttt{partially\_solved}, rigor label \texttt{proved}, novelty confidence \texttt{low}; classification \texttt{PARTIAL-PROGRESS}, registry status \texttt{partially\_solved}. Candidate: finite-subfield saturation obstruction. The description of substantial but restricted literature progress is a dated aid assessment.
\item Ulam AI / UnsolvedMath Contributors; CC BY 4.0. Machine-generated research aid, not peer reviewed; independently verify before citation.
\end{itemize}
\end{document}
